\chapter{Glossary and questions a reader may have}
\label{app:glossary}

\section{Glossary}
\begin{description}
  \item[ACE] Adaptive Continuous Entanglement simulator, built on SeQUeNCe v0.8.1.
  \item[BSM] Bell-state measurement station placed on a link between two cores; it heralds the generation of an EPR pair.
  \item[Budget] The number of communication memories per core. $1$--$4$ are modelled; above four is counterfactual ($^*$).
  \item[Common latency] Time from when a request could first be asked for until it completes, minus the 5.3\,ms lead (Chapter~\ref{ch:corrections}).
  \item[Compiler cap] The most memories per core that compiler pre-generation may hold at once (\code{budget - 1}). The memories are one shared pool: a request can use any memory that is free, including one the compiler could have held.
  \item[CGP, ACGP] Continuous and adaptive continuous generation: autonomous protocols that keep generating pairs for a neighbour drawn at random (CGP: uniformly; ACGP: with probabilities raised for neighbours seen in recently used paths), whether or not a request wants them (Section~\ref{sec:code-reading}).
  \item[ODG] On-demand generation: a pair is generated only when a request asks for it.
  \item[Congestion] Requests waiting because a core has no free communication memory.
  \item[Counterfactual] A what-if capacity beyond the modelled hardware.
  \item[Cut-off] Time after which a stored pair is discarded (Chapter~\ref{ch:cutoff}).
  \item[Dynamic] Compiler pre-generation: pairs are prepared ahead of the layer that needs them.
  \item[Endpoint demand $D_{\mathrm{end}}$] Hops touching a core that start or end a route there.
  \item[Excess] Demand above the capacity, summed over cores and layers.
  \item[HQA] The qubit-mapping algorithm (Hungarian-assignment based) that places logical qubits on cores layer by layer.
  \item[Hop] One request for an EPR pair between two adjacent cores.
  \item[Layer] One step of the mapped circuit; its hops are the moves needed to reach the next placement.
  \item[On-demand] A pair is generated only when a request asks for it.
  \item[Relay] A core strictly between a route's source and destination.
  \item[Route] A qubit's consecutive hops within one layer.
  \item[Serialization] The legacy per-core batching that stopped the budget from binding.
  \item[Torus] A mesh whose opposite edges are joined by extra links.
\end{description}

\section{Questions a reader may have}
\paragraph{Why is the latency floor 1.1--1.2\,ms?} That is the generation time of one pair once nothing waits; it is the same for every workload and does not depend on the budget.

\paragraph{Why does dynamic start slower than on-demand?} Its prepared pairs may occupy up to \code{budget - 1} memories, and a request needs a free memory of its own even when its pair is ready, so with two memories a prepared pair and the request that uses it already occupy both (Chapter~\ref{ch:why-dynamic}).

\paragraph{Why do GHZ, Bernstein--Vazirani and BMW change under dynamic between two and three memories if they have no congestion?} The structural overload is zero, so it is not congestion; the earlier budget study attributed it to planner coverage (not re-tested here).

\paragraph{Are these numbers physical?} They come from a physical simulator with ACE's parameters, but those parameters are undecided; the congestion results depend on them little, the pre-generation and fidelity results entirely (Appendix~\ref{app:params}).

\paragraph{Can the results be reproduced?} Yes: Appendix~\ref{app:repro} lists the commands, and every result folder holds the parameters, trace hashes and code commit it was produced with. The two full Experiment~1 runs (first pass and rerun) are numerically identical apart from the one column that was redefined.

\paragraph{What was not done?} Experiment~1 for fixed pre-generation and the autonomous caches; routing combined with pre-generation; joint mapping and routing; a cut-off sweep against the budget (Chapter~\ref{ch:limits}).
